\begin{trapbox}
\textbf{C016 的高频易错点及纠正}
\begin{enumerate}
  \item \textbf{左支只改 $x_I$，却不改分母。}
  错：左支写 $I(-a,cy_0/(c+ex_0))$。
  对：\[\boxed{I\left(-a,\frac{cy_0}{c-ex_0}\right)}\quad(x_0<-a).\]
  最稳妥是统一写 $I\left(a\operatorname{sgn}(x_0),ay_0/(a+|x_0|)\right)$。

  \item \textbf{将 $c/(c+ex_0)$ 错化成 $a/(a+ex_0)$。}
  因为 $e=c/a$，正确结果为
  $\displaystyle \frac{c}{c+ex_0}=\frac{a}{a+x_0}$（仅右支）。

  \item \textbf{忘记 $P$ 必须是非顶点。}
  当 $P=(\pm a,0)$，三点共线，三角形退化，通常意义下的四心不存在；
  $H_y,O'_y$ 公式也因 $y_0=0$ 无法使用。

  \item \textbf{把内心轨迹写成整条 $x=\pm a$。}
  实际还有 $0<|y_I|<b$；顶点处 $y_I=0$ 与端点 $|y_I|=b$ 均取不到。
  只说“内心在顶点切线上”是\emph{位置性质}，不是精确的\emph{轨迹集合}。

  \item \textbf{认为 $r=y_I$ 在上下半平面都成立。}
  内切圆半径 $r>0$，而 $y_I$ 可以为负；准确关系是
  $y_I=\operatorname{sgn}(y_0)r$。

  \item \textbf{把焦半径之和与差混用。}
  双曲线定义提供 $|PF_1-PF_2|=2a$，
  焦半径之和却是 $PF_1+PF_2=2e|x_0|$，并非 $2a$。
  椭圆才有 $PF_1+PF_2=2a$。

  \item \textbf{将角平分线交点 $Q$ 误当内切圆切点 $T$。}
  一般 $Q=(a^2/x_0,0)$，而 $T=(\pm a,0)$，
  它们不是同一个点。只有特殊退化极限才可能靠近。

  \item \textbf{用斜率时忘记竖直线。}
  当 $x_0=c$ 时 $PF_2$ 竖直，斜率 $y_0/(x_0-c)$ 无定义。
  求高线可直接使用
  $(x_0-c)(x+c)+y_0y=0$，不必除以零。

  \item \textbf{以为外心、垂心永远在三角形内部。}
  只有内心和重心始终位于非退化三角形内部；
  外心、垂心可以在内部、边界或外部。
  特别地，$x_0=c$ 时垂心恰好为 $F_2$。

  \item \textbf{错误判断顶点附近的极限。}
  当 $x_0\to\pm a$、$y_0\to0$ 时，由 $c^2-a^2=b^2>0$，
  $H_y\sim b^2/y_0$，$y_{O'}\sim-b^2/(2y_0)$，
  二者均可能发散（符号相反）；
  内心却趋于 $(\pm a,0)$。\textbf{极限位置不等于退化三角形有四心。}

  \item \textbf{不分两焦点的记号顺序。}
  本结论固定 $F_1=(-c,0)$、$F_2=(c,0)$。
  公式 $x_I=(PF_1-PF_2)/2$ 是在这个约定下成立的；
  若交换焦点标号，也必须同步调整带符号公式。

  \item \textbf{记住了四条公式，却没看到通用三角形结构。}
  重心靠坐标平均；外心靠中垂线；垂心靠高线或欧拉线；
  内心靠切线长与 $S=rs$。
  这些思路对焦点在 $x$ 轴上的椭圆同样适用。
\end{enumerate}
\end{trapbox}
